\documentclass[tikz,border=8pt]{standalone}
\usepackage{tikz}
\usepackage{amsmath}
\usepackage{amssymb}
\usepackage{pgfplots}
\pgfplotsset{compat=1.18}
\usepackage{ctex}
\usetikzlibrary{calc,positioning,arrows.meta,trees,matrix}

% N-Queens（LC 51/52）4×4 棋盘
% 解：皇后位置 (行,列) = (0,1),(1,3),(2,0),(3,2)（0-indexed）

\begin{document}
\begin{tikzpicture}

%% ── 标题 ───────────────────────────────────────────────────
\node[font=\large\bfseries] at (3.5, 10.0) {N 皇后（LC 52 N-Queens II）};
\node[font=\small,gray] at (3.5, 9.4)
  {$4\times4$ 棋盘，解：列 $[1,3,0,2]$，标红线为攻击范围};

%% ── 棋盘（4×4，每格 1.6cm，左下角 (0,0)） ───────────────────
\def\cellsize{1.6}

% 画格子背景（棋盘黑白交替）
\foreach \r in {0,1,2,3} {
  \foreach \c in {0,1,2,3} {
    \pgfmathparse{int(mod(\r+\c,2))}
    \ifnum\pgfmathresult=0
      \fill[gray!15] (\c*\cellsize, \r*\cellsize) rectangle (\c*\cellsize+\cellsize, \r*\cellsize+\cellsize);
    \else
      \fill[white] (\c*\cellsize, \r*\cellsize) rectangle (\c*\cellsize+\cellsize, \r*\cellsize+\cellsize);
    \fi
  }
}

% 画格子边框
\foreach \r in {0,1,2,3} {
  \foreach \c in {0,1,2,3} {
    \draw[gray!60] (\c*\cellsize, \r*\cellsize) rectangle (\c*\cellsize+\cellsize, \r*\cellsize+\cellsize);
  }
}
\draw[black,thick] (0,0) rectangle (4*\cellsize, 4*\cellsize);

% 行号（左侧）
\foreach \r in {0,1,2,3} {
  \node[font=\small,gray] at (-0.5, \r*\cellsize+0.5*\cellsize) {\r};
}
% 列号（下方）
\foreach \c in {0,1,2,3} {
  \node[font=\small,gray] at (\c*\cellsize+0.5*\cellsize, -0.5) {\c};
}

%% ── 皇后位置：(row, col) = (0,1),(1,3),(2,0),(3,2) ───────────
% 解：第 0 行放列 1，第 1 行放列 3，第 2 行放列 0，第 3 行放列 2
% TikZ 坐标：x = col*1.6 + 0.8，y = row*1.6 + 0.8

% 皇后格（黄色高亮）
\fill[yellow!40] (1*\cellsize, 0*\cellsize) rectangle (1*\cellsize+\cellsize, 0*\cellsize+\cellsize);
\fill[yellow!40] (3*\cellsize, 1*\cellsize) rectangle (3*\cellsize+\cellsize, 1*\cellsize+\cellsize);
\fill[yellow!40] (0*\cellsize, 2*\cellsize) rectangle (0*\cellsize+\cellsize, 2*\cellsize+\cellsize);
\fill[yellow!40] (2*\cellsize, 3*\cellsize) rectangle (2*\cellsize+\cellsize, 3*\cellsize+\cellsize);

% 重画边框
\foreach \r in {0,1,2,3} {
  \foreach \c in {0,1,2,3} {
    \draw[gray!60] (\c*\cellsize, \r*\cellsize) rectangle (\c*\cellsize+\cellsize, \r*\cellsize+\cellsize);
  }
}
\draw[black,thick] (0,0) rectangle (4*\cellsize, 4*\cellsize);

% 皇后符号
\node[font=\LARGE] at (1*\cellsize+0.8, 0*\cellsize+0.8) {♛};
\node[font=\LARGE] at (3*\cellsize+0.8, 1*\cellsize+0.8) {♛};
\node[font=\LARGE] at (0*\cellsize+0.8, 2*\cellsize+0.8) {♛};
\node[font=\LARGE] at (2*\cellsize+0.8, 3*\cellsize+0.8) {♛};

%% ── 攻击线（红色虚线） ───────────────────────────────────────
% 皇后 (0,1)：行 0 水平线
\draw[red!70,dashed,thick] (-0.05, 0*\cellsize+0.8) -- (4*\cellsize+0.05, 0*\cellsize+0.8);
% 皇后 (1,3)：行 1 水平线
\draw[red!70,dashed,thick] (-0.05, 1*\cellsize+0.8) -- (4*\cellsize+0.05, 1*\cellsize+0.8);
% 皇后 (2,0)：行 2 水平线
\draw[red!70,dashed,thick] (-0.05, 2*\cellsize+0.8) -- (4*\cellsize+0.05, 2*\cellsize+0.8);
% 皇后 (3,2)：行 3 水平线
\draw[red!70,dashed,thick] (-0.05, 3*\cellsize+0.8) -- (4*\cellsize+0.05, 3*\cellsize+0.8);

% 列线
\draw[red!70,dashed,thick] (1*\cellsize+0.8, -0.05) -- (1*\cellsize+0.8, 4*\cellsize+0.05);
\draw[red!70,dashed,thick] (3*\cellsize+0.8, -0.05) -- (3*\cellsize+0.8, 4*\cellsize+0.05);
\draw[red!70,dashed,thick] (0*\cellsize+0.8, -0.05) -- (0*\cellsize+0.8, 4*\cellsize+0.05);
\draw[red!70,dashed,thick] (2*\cellsize+0.8, -0.05) -- (2*\cellsize+0.8, 4*\cellsize+0.05);

% 对角线 (0,1) → ↗↙
\draw[red!50,dashed] (0.0, 1*\cellsize-0.8) -- (3.2, 4*\cellsize+0.05);
\draw[red!50,dashed] (3.2, -0.05) -- (4*\cellsize+0.05, 1*\cellsize+0.8-1*\cellsize+3.2*0);

% 对角线 (1,3) → 主对角
\draw[red!50,dashed] (3*\cellsize+0.8-1*\cellsize, 0) -- (4*\cellsize+0.05, 1.6);
\draw[red!50,dashed] (0, 1*\cellsize+0.8-3*\cellsize*0) -- (3*\cellsize+0.8, 4*\cellsize+0.05);

%% ── 三个集合标注（底部） ────────────────────────────────────
\node[draw=gray!50,fill=blue!5,rounded corners=3pt,font=\small,inner sep=7pt,align=left]
  at (3.5, -2.0)
  {\textbf{冲突检测集合（已占用）}\\[5pt]
   \textbf{cols}：$\{0,1,2,3\}$（每行放置后加入列号）\\[3pt]
   \textbf{diag1}（$row-col$）：$\{-1,\ -2,\ 2,\ 1\}$（主对角线标识）\\[3pt]
   \textbf{diag2}（$row+col$）：$\{1,\ 4,\ 2,\ 5\}$（副对角线标识）};

%% ── 合法判断说明 ────────────────────────────────────────────
\node[draw=gray!50,fill=white,rounded corners=3pt,font=\small,inner sep=6pt,align=center]
  at (3.5, -4.2)
  {放置合法条件：$col \notin \textbf{cols}$
   \quad 且 \quad $row-col \notin \textbf{diag1}$
   \quad 且 \quad $row+col \notin \textbf{diag2}$};

%% ── 标注框：坐标说明 ────────────────────────────────────────
\node[draw=gray!40,fill=white,rounded corners=2pt,font=\scriptsize,inner sep=4pt,align=left]
  at (9.0, 4.8)
  {解的列序列：$[1,3,0,2]$\\
   第 0 行放列 1\\
   第 1 行放列 3\\
   第 2 行放列 0\\
   第 3 行放列 2};

\draw[->,>=Stealth,gray!60] (8.0,4.8) -- (6.8,6.2);

\end{tikzpicture}
\end{document}
